% Fuente única de clase: Recursos/2025-11-11.pdf, páginas 1--8.
En la clase 22 estudiamos característica, grado y extensiones simples de campos. La evaluación determina el polinomio mínimo, permite extender isomorfismos y conduce al conteo de automorfismos sobre el campo base.

%%%%%%%%%%%%%%%%%%%%%%%%%%%%%%%%%%%%%%%%%%%%%%%%%%%%%%%%%%%%%%%%%%%%%%%%%%%%%%%
\section{Característica y grado de una extensión}
%%%%%%%%%%%%%%%%%%%%%%%%%%%%%%%%%%%%%%%%%%%%%%%%%%%%%%%%%%%%%%%%%%%%%%%%%%%%%%%
%%%%%%%%%%%%%%%%%%%%%%%%%%%%%%%%%%%%%%%%%%%%%%%%%%%%%%%%%%%%%%%%%%%%%%%%%%%%%%%
\subsection{Morfismos y campos primos}
%%%%%%%%%%%%%%%%%%%%%%%%%%%%%%%%%%%%%%%%%%%%%%%%%%%%%%%%%%%%%%%%%%%%%%%%%%%%%%%
% P01--P02.
Todo homomorfismo de campos $i:k\to K$ es inyectivo, como se comprobó en la clase 18: su núcleo es un ideal de $k$ y no contiene a $1$. Así, $i$ identifica $k$ con un subcampo de $K$ y define una \emph{extensión de campos}. Escribiremos $k\subseteq K$, recordando que la inclusión representa un homomorfismo concreto; puede haber más de una forma de incluir el mismo campo abstracto. El símbolo $K/k$, cuando se usa para una extensión, no denota un cociente de anillos.

El único homomorfismo unitario $j:\mathbb Z\to k$ viene dado por $j(n)=n1_k$. La característica $\operatorname{char}k$ es el generador no negativo de $\ker j$. Como la imagen es un subanillo del dominio $k$, ese núcleo es primo, luego es $(0)$ o $(p)$ para un primo positivo $p$. Así, la característica de un campo es cero o un primo. Por ejemplo,
\[
\operatorname{char}\mathbb Q=0,\qquad
\operatorname{char}\mathbb F_p=p,\qquad
\mathbb F_p=\mathbb Z/p\mathbb Z.
\]
La característica se conserva al extender campos. En efecto, en el cuadrado de la página 1,
\[
\begin{tikzcd}[column sep=large,row sep=large]
\mathbb Z\arrow[r,"\id"]\arrow[d,"j"']&\mathbb Z\arrow[d,"j'"]\\
k\arrow[r,hook,"i"']&K,
\end{tikzcd}
\]
se tiene $j'=i\circ j$; la inyectividad de $i$ implica $\ker j'=\ker j$.

\begin{prop}\label{prop:C22-campo-primo}
Todo campo de característica cero contiene una única copia de $\mathbb Q$ compatible con la unidad; todo campo de característica $p$ contiene una única copia de $\mathbb F_p$. Estos campos son los objetos iniciales en las categorías de campos de su respectiva característica.
\end{prop}
\begin{proof}
Si $\operatorname{char}F=0$, el homomorfismo $\mathbb Z\to F$ es inyectivo. Por la propiedad universal del campo de fracciones se extiende de manera única a $\mathbb Q\to F$, mediante $a/b\mapsto j(a)j(b)^{-1}$:
\[
\begin{tikzcd}[column sep=large,row sep=large]
\mathbb Q\arrow[r,hook,"\widetilde j",destacado]&F\\
\mathbb Z\arrow[u,hook]\arrow[ur,"j"']&
\end{tikzcd}
\]
Si $\operatorname{char}F=p$, el primer isomorfismo induce $\mathbb Z/(p)\cong\im j\subseteq F$. La aplicación está determinada por $1\mapsto1$, así que es única. En ambos casos es inyectiva por ser un homomorfismo de campos.
\end{proof}
Escribimos $\mathbf{Fld}_c$ para la categoría de campos de característica $c$, donde $c=0$ o $c=p$ primo. La proposición expresa que $\mathbb Q$ es inicial en $\mathbf{Fld}_0$ y $\mathbb F_p$ en $\mathbf{Fld}_p$. La copia correspondiente dentro de $F$ se llama su \emph{campo primo}.

%%%%%%%%%%%%%%%%%%%%%%%%%%%%%%%%%%%%%%%%%%%%%%%%%%%%%%%%%%%%%%%%%%%%%%%%%%%%%%%
\subsection{Dimensión sobre el campo base}
%%%%%%%%%%%%%%%%%%%%%%%%%%%%%%%%%%%%%%%%%%%%%%%%%%%%%%%%%%%%%%%%%%%%%%%%%%%%%%%
% P02--P03.
\begin{defi}
Una extensión $k\subseteq F$ es \emph{finita de grado $n$} si $F$, como espacio vectorial sobre $k$, tiene dimensión finita $n$. Escribimos $[F:k]=n$. Si la dimensión es infinita, la extensión es \emph{infinita} y en esta clase escribimos $[F:k]=\infty$.
\end{defi}
El grado mide la dimensión, no el número de elementos del campo. La estructura vectorial proviene de la multiplicación de $F$ restringida a escalares de $k$.

\begin{prop}\label{prop:C22-grado-cociente}
Si $f\in k[t]$ es irreducible y $n=\deg f\geq1$, entonces $F=k[t]/(f)$ es un campo y la inclusión $k\to F$ define una extensión de grado $n$. Si $\alpha=t+(f)$, una base es $1,\alpha,\ldots,\alpha^{n-1}$.
\end{prop}
\begin{proof}
El anillo $k[t]$ es euclidiano, luego DIP. Por ello el ideal generado por el irreducible $f$ es maximal y su cociente es campo. Los elementos constantes se incluyen inyectivamente: un múltiplo no nulo de $f$ no puede ser una constante. Al dividir por $f$, cada clase tiene un representante único de grado menor que $n$, por lo que se expresa de manera única como combinación de las potencias indicadas. Esto prueba tanto generación como independencia.
\end{proof}
En particular, $\mathbb C\cong\mathbb R[t]/(t^2+1)$ da $[\mathbb C:\mathbb R]=2$. La construcción produce extensiones concretas, pero no afirma que toda extensión sea un cociente de esta forma.

%%%%%%%%%%%%%%%%%%%%%%%%%%%%%%%%%%%%%%%%%%%%%%%%%%%%%%%%%%%%%%%%%%%%%%%%%%%%%%%
\section{Extensiones simples y polinomios mínimos}
%%%%%%%%%%%%%%%%%%%%%%%%%%%%%%%%%%%%%%%%%%%%%%%%%%%%%%%%%%%%%%%%%%%%%%%%%%%%%%%
%%%%%%%%%%%%%%%%%%%%%%%%%%%%%%%%%%%%%%%%%%%%%%%%%%%%%%%%%%%%%%%%%%%%%%%%%%%%%%%
\subsection{El campo generado por un elemento}
%%%%%%%%%%%%%%%%%%%%%%%%%%%%%%%%%%%%%%%%%%%%%%%%%%%%%%%%%%%%%%%%%%%%%%%%%%%%%%%
% P03--P04.
Para $\alpha\in F$ y $k\subseteq F$, definimos $k(\alpha)$ como el menor subcampo de $F$ que contiene a $k$ y a $\alpha$. Existe porque la intersección de todos los subcampos con esas propiedades es un subcampo; la familia no es vacía, pues incluye a $F$.
\begin{defi}
La extensión $k\subseteq F$ es \emph{simple} si $F=k(\alpha)$ para algún $\alpha\in F$.
\end{defi}
En el cociente de la proposición anterior, todo subcampo que contenga a $k$ y a $\alpha=t+(f)$ contiene todas las expresiones polinómicas en $\alpha$, que son todos los elementos del cociente. Por tanto $F=k(\alpha)$. Además,
\[
f(\alpha)=f\bigl(t+(f)\bigr)=f(t)+(f)=0,
\]
y la inclusión del campo base es la composición $k\hookrightarrow k[t]\twoheadrightarrow k[t]/(f)$.

\begin{nota}[Sentido de la inclusión]
En la parte manuscrita superior de la página 4, después de observar que un subcampo $F'$ contiene todas las clases polinómicas, aparece la inclusión $F'\subseteq k(\alpha)$. Lo que prueba ese argumento es $F\subseteq F'$. Como ya $F'\subseteq F$, resulta $F'=F$ y por tanto $k(\alpha)=F$.
\end{nota}

\begin{samepage}
%%%%%%%%%%%%%%%%%%%%%%%%%%%%%%%%%%%%%%%%%%%%%%%%%%%%%%%%%%%%%%%%%%%%%%%%%%%%%%%
\subsection{La dicotomía de la evaluación}
%%%%%%%%%%%%%%%%%%%%%%%%%%%%%%%%%%%%%%%%%%%%%%%%%%%%%%%%%%%%%%%%%%%%%%%%%%%%%%%
% P04--P05.
\begin{prop}\label{prop:C22-simple}
Sea $k\subseteq k(\alpha)$ una extensión simple y sea
\[
\epsilon:k[t]\to k(\alpha),\qquad g\longmapsto g(\alpha).
\]
Si $\epsilon$ es inyectiva, la extensión es infinita y $k(\alpha)\cong k(t)$ sobre $k$, donde $k(t)$ es el campo de funciones racionales. Si $\epsilon$ no es inyectiva, existe un único polinomio mónico irreducible no constante $p\in k[t]$ que genera $\ker\epsilon$, y
\[
k(\alpha)\cong k[t]/(p),\qquad [k(\alpha):k]=\deg p.
\]
Este $p$ es el polinomio mónico de menor grado que se anula en $\alpha$. En particular, la evaluación es inyectiva si y sólo si la extensión es infinita.
\end{prop}
\end{samepage}
\begin{proof}
Por el primer isomorfismo,
\[
k[t]/\ker\epsilon\cong\im\epsilon\subseteq k(\alpha).
\]
La imagen es un dominio, por lo que $\ker\epsilon$ es un ideal primo.

Si el núcleo es cero, la propiedad universal de fracciones extiende $\epsilon$ a una inclusión
\[
\begin{tikzcd}[column sep=large,row sep=large]
k(t)\arrow[r,hook,"\widetilde\epsilon",destacado]&k(\alpha)\\
k[t]\arrow[u,hook]\arrow[ur,hook,"\epsilon"']&
\end{tikzcd}
\qquad
\widetilde\epsilon\left(\frac{g(t)}{h(t)}\right)
=\frac{g(\alpha)}{h(\alpha)}.
\]
El denominador es no nulo por inyectividad. La imagen de $\widetilde\epsilon$ es un subcampo que contiene a $k$ y a $\alpha$, de modo que coincide con $k(\alpha)$ por su minimalidad. Además, $1,t,t^2,\ldots$ es independiente sobre $k$ en $k[t]$; su imagen $1,\alpha,\alpha^2,\ldots$ lo sigue siendo. La extensión tiene dimensión infinita.

Si el núcleo es no nulo, por ser $k[t]$ DIP hay un único generador mónico $p$. Es irreducible porque genera un ideal primo no nulo, y no es constante porque el núcleo es propio. El cociente $k[t]/(p)$ es campo; por tanto $\im\epsilon$ es un subcampo de $k(\alpha)$ que contiene a $k$ y a $\alpha$, y nuevamente es todo $k(\alpha)$. La \cref{prop:C22-grado-cociente} da el grado y la base de potencias. Como todo polinomio que se anula en $\alpha$ es múltiplo de $p$, ninguno no nulo puede tener grado menor. Si otro mónico del mismo grado se anula, el cociente de la división es $1$ y el resto cero, así que coincide con $p$.
\end{proof}
El grado que se calcula al final de la página 5 debe escribirse $[k(\alpha):k]$, como en el enunciado impreso. Allí la última línea manuscrita pone $k[t]$ en el segundo lugar, pero éste no es el campo base de esa dimensión.

\begin{defi}
En el caso finito, el polinomio $p$ anterior es el \emph{polinomio mínimo de $\alpha$ sobre $k$}. Depende del campo base y es mónico por definición.
\end{defi}

%%%%%%%%%%%%%%%%%%%%%%%%%%%%%%%%%%%%%%%%%%%%%%%%%%%%%%%%%%%%%%%%%%%%%%%%%%%%%%%
\subsection{Dos inclusiones para una raíz cuadrada}
%%%%%%%%%%%%%%%%%%%%%%%%%%%%%%%%%%%%%%%%%%%%%%%%%%%%%%%%%%%%%%%%%%%%%%%%%%%%%%%
% P06.
Tomemos $\mathbb Q\subseteq\mathbb R$. El polinomio $t^2-2$ es irreducible sobre $\mathbb Q$: una raíz racional tendría denominador uno y sería un entero que divide a $2$, y ninguno de $\pm1,\pm2$ lo anula. Las dos raíces reales producen dos homomorfismos distintos
\[
\overline\epsilon_\pm:\mathbb Q[t]/(t^2-2)\hookrightarrow\mathbb R,
\qquad t+(t^2-2)\longmapsto\pm\sqrt2.
\]
Son las inclusiones estudiadas en la clase 19. La imagen de la positiva es $\mathbb Q(\sqrt2)$ y la negativa tiene imagen $\mathbb Q(-\sqrt2)$; en este ejemplo son el mismo subcampo. Sus elementos son
\[
\mathbb Q(\sqrt2)=\set{a+b\sqrt2\given a,b\in\mathbb Q}.
\]
La evaluación y el paso al cociente conservan el diagrama de la página 6:
\[
\begin{tikzcd}[column sep=large,row sep=large]
\mathbb Q[t]\arrow[r,"\epsilon_\pm"]\arrow[d,two heads,"\pi"']&\mathbb R\\
\mathbb Q[t]/(t^2-2)\arrow[ur,hook,"\overline\epsilon_\pm"',destacado]&
\end{tikzcd}
\]
La inclusión $(t^2-2)\subseteq\ker\epsilon_\pm$ permite inducir el mapa; como el polinomio es mínimo, ambos ideales son iguales. Fijar una raíz fija el mediador de manera única, pero cambiarla puede cambiar el homomorfismo. Para raíces distintas de un polinomio irreducible, los campos generados dentro de un campo ambiente son isomorfos sobre la base; no se deduce que sean el mismo subcampo.

%%%%%%%%%%%%%%%%%%%%%%%%%%%%%%%%%%%%%%%%%%%%%%%%%%%%%%%%%%%%%%%%%%%%%%%%%%%%%%%
\section{Extensión de isomorfismos y automorfismos}
%%%%%%%%%%%%%%%%%%%%%%%%%%%%%%%%%%%%%%%%%%%%%%%%%%%%%%%%%%%%%%%%%%%%%%%%%%%%%%%
%%%%%%%%%%%%%%%%%%%%%%%%%%%%%%%%%%%%%%%%%%%%%%%%%%%%%%%%%%%%%%%%%%%%%%%%%%%%%%%
\subsection{Transportar el polinomio mínimo}
%%%%%%%%%%%%%%%%%%%%%%%%%%%%%%%%%%%%%%%%%%%%%%%%%%%%%%%%%%%%%%%%%%%%%%%%%%%%%%%
% P07--P08.
\begin{prop}\label{prop:C22-extender-iso}
Sean $F_1=k_1(\alpha_1)$ y $F_2=k_2(\alpha_2)$ extensiones simples finitas, con polinomios mínimos $p_1$ y $p_2$. Si $i:k_1\to k_2$ es un isomorfismo que transforma los coeficientes de $p_1$ en los de $p_2$, existe un único isomorfismo $j:F_1\to F_2$ que extiende a $i$ y satisface $j(\alpha_1)=\alpha_2$.
\end{prop}
\begin{proof}
La transformación de coeficientes define un isomorfismo de anillos $\widehat i:k_1[t]\to k_2[t]$, que fija $t$. Por hipótesis lleva $(p_1)$ sobre $(p_2)$. La composición con la proyección a $k_2[t]/(p_2)$ es sobreyectiva y tiene núcleo $(p_1)$, de modo que induce un isomorfismo de cocientes. Identificándolos con $F_1$ y $F_2$ obtenemos el cuadrado de la página 7:
\[
\begin{tikzcd}[column sep=large,row sep=large]
k_1[t]/(p_1)\arrow[r,"\overline i","\sim"']\arrow[d,"\sim"']&k_2[t]/(p_2)\arrow[d,"\sim"]\\
k_1(\alpha_1)\arrow[r,"j","\sim"',destacado]&k_2(\alpha_2).
\end{tikzcd}
\]
El mapa así construido fija los datos requeridos. Para la unicidad, sea $n=\deg p_1=\deg p_2$. Todo $x\in F_1$ tiene una expresión única
\[
x=c_0+c_1\alpha_1+\cdots+c_{n-1}\alpha_1^{n-1}.
\]
Un homomorfismo que extienda a $i$ y envíe $\alpha_1$ a $\alpha_2$ debe satisfacer
\[
j(x)=i(c_0)+i(c_1)\alpha_2+\cdots+i(c_{n-1})\alpha_2^{n-1},
\]
lo que determina toda la aplicación.
\end{proof}
Decimos que $j$ \emph{extiende} o \emph{levanta} a $i$. El diagrama de extensiones de la página 8 es
\[
\begin{tikzcd}[column sep=large,row sep=large]
k_1(\alpha_1)\arrow[r,"j","\sim"',destacado]&k_2(\alpha_2)\\
k_1\arrow[u,hook]\arrow[r,"i","\sim"']&k_2\arrow[u,hook].
\end{tikzcd}
\]

%%%%%%%%%%%%%%%%%%%%%%%%%%%%%%%%%%%%%%%%%%%%%%%%%%%%%%%%%%%%%%%%%%%%%%%%%%%%%%%
\subsection{Contar automorfismos mediante raíces}
%%%%%%%%%%%%%%%%%%%%%%%%%%%%%%%%%%%%%%%%%%%%%%%%%%%%%%%%%%%%%%%%%%%%%%%%%%%%%%%
% P08.
Nos interesan en particular los isomorfismos de un campo consigo mismo que fijan un subcampo dado. Para $k\subseteq F$ escribimos
\[
\Aut_k(F)=\set{\sigma\in\Aut(F)\given \sigma(c)=c\text{ para todo }c\in k}.
\]
Son los automorfismos de la extensión como objeto en la categoría de campos sobre $k$, que aquí denotamos $k\text{-}\mathbf{Fld}$ para distinguirla de $\mathbf{Fld}_c$, donde el subíndice indica una característica.

\begin{coro}\label{coro:C22-automorfismos}
Si $F=k(\alpha)$ es una extensión simple finita y $p$ es el polinomio mínimo de $\alpha$, el número $|\Aut_k(F)|$ es el número de raíces distintas de $p$ en $F$. En particular,
\[
|\Aut_k(F)|\leq[F:k],
\]
con igualdad si y sólo si $p$ se factoriza en $F[t]$ como producto de factores lineales distintos.
\end{coro}
\begin{proof}
Todo $\sigma\in\Aut_k(F)$ queda determinado por $\sigma(\alpha)$, pues cada elemento de $F$ es una expresión polinómica en $\alpha$ con coeficientes en $k$. Además,
\[
p\bigl(\sigma(\alpha)\bigr)=\sigma\bigl(p(\alpha)\bigr)=\sigma(0)=0.
\]
Por tanto $\sigma\mapsto\sigma(\alpha)$ es una aplicación inyectiva hacia el conjunto de raíces de $p$ en $F$.

Si $\beta\in F$ es otra raíz, su polinomio mínimo sobre $k$ es también $p$: su ideal de evaluación contiene al irreducible $p$, y el generador mónico mínimo debe ser ese mismo polinomio. La \cref{prop:C22-extender-iso} da un isomorfismo $k(\alpha)\to k(\beta)$ que fija $k$ y envía $\alpha$ a $\beta$. Como $k(\beta)\subseteq F$ y ambos tienen dimensión $\deg p$ sobre $k$, se tiene $k(\beta)=F$. Así, el isomorfismo es un automorfismo de $F$, y toda raíz se obtiene de esa manera. Se ha explicitado este paso de sobreyectividad de la prueba impresa.

Finalmente, $p$ tiene a lo sumo $\deg p=[F:k]$ raíces distintas. Alcanza esa cantidad exactamente cuando se descompone por completo en factores lineales distintos, por la división sucesiva entre $t-\beta$ y la igualdad de grados.
\end{proof}

\begin{obs}[Orientación bibliográfica]
Los fragmentos de \citet{aluffi2009} en las ocho páginas pertenecen a la apertura de VII y a VII.1.1--2: definición VII.1.1 de grado, definición VII.1.2 de extensión simple, proposición VII.1.3, ejemplo VII.1.4, proposición VII.1.5 y corolario VII.1.7. Se desarrollan la característica, los objetos iniciales, el grado del cociente, la evaluación en sus dos casos, el transporte de isomorfismos y el conteo de automorfismos. La remisión inicial a V.4.17 no reproduce su enunciado; se justifican las propiedades de característica efectivamente usadas, sin reconstruir un ejercicio no visible. La remisión a V.5.7 corresponde a la adjunción ya estudiada en la clase 19. Las precisiones a las páginas 4--5 y el paso dimensional del corolario se indican en el texto. No se consultó el libro completo.
\end{obs}
